% Lösungen Fokus Nullstellen – Einheit 4
\einheitenkopf[][4 Nullstellen]{Einheit 4 · Nullstellen und Schnittpunkte berechnen}

\erg{6}{a) $0 = (x-5)^2 - 4$ \quad b) $0 = x^2 + 3x - 10$ \quad c) $0 = -(x+1)^2 + 9$ \quad d) $0 = x^2 - 36$ \quad e) $0 = 3x^2 - 12x + 9$}
\erg{7}{a) $x_1 = 3$, $x_2 = -1$ \quad b) $x_1 = 5$, $x_2 = 3$ \quad c) $x_1 = 5$, $x_2 = -1$ \quad d) $x_1 = 9$, $x_2 = 1$}
\erg{8}{a) $x_1 = -1$, $x_2 = -3$ \quad b) $x_1 = 1$, $x_2 = -9$ \quad c) $16 = (x-1)^2$; $x_1 = 5$, $x_2 = -3$ \quad d) $4 = (x+3)^2$; $x_1 = -1$, $x_2 = -5$}
\erg{9}{a) $(x-3)^2 = 4$; $x_1 = 5$, $x_2 = 1$ \quad b) $(x+1)^2 = 9$; $x_1 = 2$, $x_2 = -4$ \quad c) $(x+3)^2 = 16$; $x_1 = 1$, $x_2 = -7$}
\erg{10}{a) $x = 3 \pm \sqrt{6}$; $x_1 \approx 5{,}45$, $x_2 \approx 0{,}55$ \quad b) $x = -1 \pm \sqrt{7}$; $x_1 \approx 1{,}65$, $x_2 \approx -3{,}65$ \quad c) $0 = (x+4)^2 - 3$; $x = -4 \pm \sqrt{3}$; $x_1 \approx -2{,}27$, $x_2 \approx -5{,}73$}
\erg{11}{a) Fehler: nur die positive Wurzel genommen, die zweite Lösung fehlt. Richtig: $x - 3 = 5$ oder $x - 3 = -5$; $x_1 = 8$, $x_2 = -2$ \quad b) Fehler: Vorzeichen aus der Klammer übernommen; $x + 2 = 3$ oder $x + 2 = -3$; $x_1 = 1$, $x_2 = -5$}
\erg{12}{a) keine: Scheitel $(4 \mid 3)$ über der x-Achse, nach oben geöffnet \quad b) zwei: Scheitel $(-2 \mid -5)$ unter der x-Achse, nach oben geöffnet \quad c) eine: Scheitel $(6 \mid 0)$ liegt auf der x-Achse \quad d) zwei: Scheitel $(1 \mid 2)$ über der x-Achse, nach unten geöffnet \quad e) keine: Scheitel $(-3 \mid -4)$ unter der x-Achse, nach unten geöffnet}
\erg{13}{a) $x = 3 \pm 1$; $x_1 = 4$, $x_2 = 2$ \quad b) $x = -2 \pm 1$; $x_1 = -1$, $x_2 = -3$ \quad c) $x = 4 \pm 2$; $x_1 = 6$, $x_2 = 2$ \quad d) $x = -5 \pm 4$; $x_1 = -1$, $x_2 = -9$}
\erg{14}{a) $x = 1 \pm 3$; $x_1 = 4$, $x_2 = -2$ \quad b) $x = -3 \pm 4$; $x_1 = 1$, $x_2 = -7$ \quad c) $x = -0{,}5 \pm 2{,}5$; $x_1 = 2$, $x_2 = -3$ \quad d) $x = 1{,}5 \pm 3{,}5$; $x_1 = 5$, $x_2 = -2$}
\erg{15}{a) $x = 5 \pm \sqrt{0}$; nur $x = 5$ \quad b) $x = -4 \pm \sqrt{0}$; nur $x = -4$}
\erg{16}{a) $x = -1 \pm \sqrt{-4}$; keine Nullstelle \quad b) $x = 2 \pm \sqrt{-3}$; keine Nullstelle}
\erg{17}{a) $0 = x^2 + 2x - 15$; $x_1 = 3$, $x_2 = -5$ \quad b) $0 = x^2 + 4x - 12$; $x_1 = 2$, $x_2 = -6$ \quad c) $0 = x^2 + 2x - 8$; $x_1 = 2$, $x_2 = -4$ \quad d) $0 = x^2 + 3x + 2$; $x = -1{,}5 \pm 0{,}5$; $x_1 = -1$, $x_2 = -2$}
\erg{18}{a) $x = 2 \pm \sqrt{3}$; $x_1 \approx 3{,}73$, $x_2 \approx 0{,}27$ \quad b) $x = -1 \pm \sqrt{5}$; $x_1 \approx 1{,}24$, $x_2 \approx -3{,}24$ \quad c) $x = 5 \pm \sqrt{11}$; $x_1 \approx 8{,}32$, $x_2 \approx 1{,}68$}
\erg{19}{a) Fehler: Vorzeichen von $-\frac{p}{2}$; richtig $x = 4 \pm 3$; $x_1 = 7$, $x_2 = 1$ \quad b) Fehler: nicht durch 2 geteilt; $0 = x^2 + 6x + 5$, $x = -3 \pm 2$; $x_1 = -1$, $x_2 = -5$ \quad c) Fehler: Vorzeichen von q unter der Wurzel; $4 + 21 = 25$, $x = -2 \pm 5$; $x_1 = 3$, $x_2 = -7$}
